% Part: incompleteness
% Chapter: representability-in-q
% Section: c-representable

\documentclass[../../../include/open-logic-section]{subfiles}

\begin{document}

\olfileid{inc}{req}{cre}
\olsection{$C$ मधील फलने~$\Th{Q}$ मध्ये निरूपणीय असतात}

$C$ मधील प्रत्येक फलन~$\Th{Q}$ मध्ये निरूपणीय आहे, हे दाखवायचे
आहे. अखेरीस $C$ मधील प्रत्येक $k$-आदानी फलन
$f(x_0,\dots,x_{k-1})$ ला त्याचे निरूपण करणारे सूत्र
$!A_f(x_0,\dots,x_{k-1},y)$ कसे द्यायचे, हे दाखवावे लागेल.

\begin{lem}
$n$ आणि $m$ या नैसर्गिक संख्या असू द्या. $n \neq m$ असेल,
तर $Q \Proves \num n \neq \num m$.
\end{lem}

\begin{proof}
$n$ वर विगमन वापरून प्रत्येक $m$ साठी $n \neq m$ असेल, तर
$Q \Proves \eq/[\num n][\num m]$, हे दाखवू.

आधार टप्प्यात $n = 0$. जर $m$ हा $0$ च्या बरोबर नसेल,
तर एखाद्या नैसर्गिक संख्या $k$ साठी $m = k +
1$. आपल्याकडे $\lforall[x][\eq/[0][x']]$ असे सांगणारे स्वयंसिद्धक
आहे. संख्यापकाच्या स्वयंसिद्धकाने $x$ च्या जागी
$\num k$ ठेवून $\eq/[0][\num k']$ हा निष्कर्ष मिळतो. पण
$\num k'$ म्हणजेच $\num m$.

विगमन टप्प्यात विधान $n$ साठी खरे आहे असे मानून $n+1$
पाहू. $m$ ही कोणतीही नैसर्गिक संख्या असू द्या. दोन शक्यता आहेत:
$m = 0$, किंवा एखाद्या $k$ साठी $m = k+1$. पहिली बाब
वरप्रमाणे हाताळता येते. दुसऱ्या बाबतीत $n+1 \neq
k+1$ असे समजा. मग $n \neq k$. $n$ साठीच्या विगमन गृहीतकाने
$\Th{Q} \Proves \eq/[\num n][\num k]$. आपल्याकडे
$\lforall[x][\lforall[y][\eq[x'][y'] \lif \eq[x][y]]]$
असे सांगणारे स्वयंसिद्धक आहे. संख्यापकाच्या स्वयंसिद्धकाने
$\eq[\num n'][\num k'] \lif \eq[\num n][\num
  k]$ मिळते. विधानीय तर्कशास्त्र वापरून $\Th{Q}$ मध्ये
$\eq/[\num n][\num k] \lif \eq/[\num n'][\num k']$
हे निष्पन्न होते. विध्यात्मक अनुमानाने
$\eq/[\num n'][\num k']$ मिळते; आणि $\num k'$ म्हणजे
$\num m$ असल्यामुळे हेच अपेक्षित होते.
\end{proof}

या पूर्वप्रमेयाचा आशय मर्यादित आहे: वेगवेगळे संख्यांक
वेगवेगळ्या वस्तू दर्शवतात, हे $\Th{Q}$ सिद्ध करू शकते. उदाहरणार्थ,
$\Th{Q}$ हे $0'' \neq 0'''$ सिद्ध करते. पण हे सर्वसाधारणपणे
दाखवण्यासाठी थोडी काळजी घ्यावी लागते. येथे वापरलेले विगमन
$\Th{Q}$ च्या \emph{बाहेर} आहे, हेही लक्षात घ्या.

शून्य, उत्तरवर्ती, बेरीज, गुणाकार, समतेचे जातिबोधक फल
आणि प्रक्षेपण यांचे निरूपण करता येईल. प्रत्येक बाबतीत योग्य
निरूपक सूत्र सहज दिसते. उदाहरणार्थ, शून्याचे निरूपण
\[
y = 0,
\]
उत्तरवर्तीचे निरूपण
\[
x_0' = y,
\]
आणि बेरीजेचे निरूपण
\[
x_0 + x_1 = y.
\]
खरे काम म्हणजे संबंधित विधाने $\Th{Q}$ सिद्ध करू शकते, हे
दाखवणे. उदाहरणार्थ, वरील सूत्र बेरीजेचे निरूपण करते असे
म्हणण्यासाठी नैसर्गिक संख्यांच्या प्रत्येक जोडी $m$ आणि $n$
साठी $\Th{Q}$ पुढील सिद्ध करते, हे दाखवावे लागते:
\[
\eq[\num n + \num m][\num {n+m}]
\]
आणि
\[
\lforall[y][(\eq[(\num n + \num m)][y] \lif \eq[y][\num {n+m}])].
\]

संयोजनाचे काय? $h$ पुढीलप्रमाणे परिभाषित आहे असे माना:
\[
h(x_0,\dots,x_{l-1}) = f(g_0(x_0,\dots,x_{l-1}), \dots,
g_{k-1}(x_0,\dots,x_{l-1})).
\]
येथे $f,g_0,\dots,g_{k-1}$ यांचे अनुक्रमे निरूपण करणारी
सूत्रे $!A_f, !A_{g_0}, \dots,
!A_{g_{k-1}}$ आपल्याला आधीच मिळाली आहेत. आता $h$ चे निरूपण
करणारे सूत्र $!A_h$ परिभाषित करता येते. त्यासाठी
$!A_h(x_0,\dots,x_{l-1},y)$ पुढील असे घेऊ:
\begin{multline*}
  \lexists[z_0\dots][\lexists[z_{k-1}][(!A_{g_0}(x_0,\dots,x_{l-1},z_0) \land
  \dots \land !A_{g_{k-1}}(x_0,\dots,x_{l-1},z_{k-1}) \land]]\\
  !A_f(z_0,\dots,z_{k-1},y)).
\end{multline*}
%READERNOTE{OLINC-328: गोठवलेल्या स्रोताच्या या संयोजन-सूत्रात संख्यापकांची कंसबांधणी आणि संयुक्त-विधानाचा शेवट विस्कळीत आहे; अचूक स्रोत-सूत्र येथे जतन केले आहे.}

शेवटी अपरिबद्ध शोध पाहू. $g(x,\vec z)$ नियमित आहे आणि
$\Th{Q}$ मध्ये, समजा $!A_g(x,\vec
z,y)$ या सूत्राने, निरूपणीय आहे असे माना. $f$ ची व्याख्या
$f(\vec z) = \mu x \; g(x,\vec z)$ अशी असू द्या. $f$ चे
निरूपण करणारे सूत्र $!A_f(\vec z,y)$ शोधायचे आहे. नैसर्गिक
निवड अशी:
\[
!A_f(\vec z,y) \ident !A_g(y,\vec z,0) \land \lforall[w][(w < y
\lif \lnot !A_g(w,\vec z,0))].
\]

आणखी काही पूर्वप्रमेये वापरून हे सूत्र योग्य आहे असे दाखवता येते; उदाहरणार्थ:

\begin{lem}
  प्रत्येक चल $x$ आणि प्रत्येक नैसर्गिक संख्या $n$ साठी $\Th{Q}$
  हे $\eq[(x'
  + \num n)][(x + \num n)']$ सिद्ध करते.
\end{lem}

पुन्हा लक्षात घ्या: यावरून $\Th{Q}$ हे
$\lforall[x][\lforall[y][(x' + y) = (x +y)']]$
सिद्ध करते असे म्हणता येत नाही; ते विधान खोटे आहे.

\begin{proof}
नेहमीप्रमाणे $n$ वर विगमन वापरू. आधार टप्प्यात $n =
0$ असेल, तर $\Th{Q}$ हे $\eq[(x'+0)][(x+0)']$
सिद्ध करते, हे दाखवायचे आहे. पुढील विधाने मिळतात:
\begin{eqnarray*}
& & \eq[(x' + 0)][x'] \quad \text{स्वयंसिद्धक 4 वरून} \\
& & \eq[(x + 0)][x] \quad \text{स्वयंसिद्धक 4 वरून} \\
& & \eq[(x+0)'][x'] \quad \text{ओळ 2 वरून} \\
& & \eq[(x' + 0)][(x + 0)'] \quad \text{ओळी 1 आणि 3 वरून}
\end{eqnarray*}
विगमन टप्प्यात $\Th{Q}$ मध्ये $\eq[(x' +
  \num n)][(x + \num n)']$ आधी सिद्ध झाले आहे असे
माना. $\num{n+1}$ म्हणजे $\num
n'$ असल्यामुळे $\Th{Q}$ हे
$\eq[(x' + \num n')][(x + \num
n')']$ सिद्ध करते, हे दाखवायचे आहे. पुढील समताशृंखला
$\Th{Q}$ मध्ये सिद्ध करता येते:
\begin{eqnarray*}
(x' + \num n') & \eq & (x' + \num n)' \quad \text{स्वयंसिद्धक 5} \\
& \eq & (x + \num n')' \quad \text{विगमन गृहीतकावरून}
\end{eqnarray*}
\end{proof}

\begin{lem}
\begin{enumerate}
\item $\Th{Q}$ हे $\lnot (x < \num 0)$ सिद्ध करते.
\item प्रत्येक नैसर्गिक संख्या $n$ साठी $\Th{Q}$ पुढील सिद्ध करते:
\[
x < \num {n+1} \lif (\eq[x][\num 0] \lor \dots \lor \eq[x][\num n]).
\]
\end{enumerate}
\end{lem}

\begin{proof}
१ आणि २ चा काही भाग अनौपचारिकपणे पाहू; म्हणजे औपचारिक
!!{derivation} कसे रचावे यासाठी फक्त संकेत देऊ.

भाग १ साठी $<$ च्या व्याख्येनुसार $\Th{Q}$ मध्ये
$\lnot \lexists[y][\eq[(y'+x)][0]]$ सिद्ध करायचे आहे. प्रथम-क्रम
तर्कशास्त्राच्या स्वयंसिद्धकांनी आणि नियमांनी हे
$\lforall[y][\eq/[(y' +
    x)][0]]$ याच्याशी समतुल्य आहे. कल्पना अशी: $\eq[(y'+x)][0]$
असे माना. $x$ हा 0 असेल, तर $\eq[(y' + 0)][0]$.
पण $\Th{Q}$ च्या स्वयंसिद्धक 4 वरून
$\eq[(y' + 0)][y']$, आणि स्वयंसिद्धक 2 वरून
$\eq/[y'][0]$; हा व्याघात. म्हणून
$\lforall[y][\eq/[(y' +x)][0]]$. $x$ हा $0$ नसेल,
तर स्वयंसिद्धक 3 वरून असा $z$ असतो की
$\eq[x][z']$. मग $\eq[(y' + z')][0]$.
स्वयंसिद्धक~5 वरून $\eq[(y' + z)'][0]$, हेही
स्वयंसिद्धक~2 शी विसंगत आहे.

भाग २ साठी $n$ वर विगमन वापरा. प्रथम $n =0$
हा आधार टप्पा पाहू. त्या वेळी
$x < \num 1 \lif \eq[x][\num
  0]$ दाखवायचे आहे. $x < \num 1$ असे माना.
$<$ च्या परिभाषक स्वयंसिद्धकाने
$\lexists[y][\eq[(y'+x)][0']]$ मिळते. $y$ कडे हा
गुणधर्म आहे असे माना; म्हणजे $\eq[y'+x][0']$.
%READERNOTE{OLINC-329: स्रोताचा या आधार टप्प्यातील पुढील युक्तिवाद फक्त डावीकडून उजवीकडे दिशा दाखवतो; दाखवायच्या द्विपक्षीय विधानाची उलट दिशा येथे सिद्ध केलेली नाही.}

$\eq[x][0]$ दाखवायचे आहे. स्वयंसिद्धक 3 वरून
$x$ हा 0 नसेल, तर एखाद्या $z$ साठी तो $z'$ च्या
बरोबर असतो. मग $\eq[(y' + z')][0']$.
$\Th{Q}$ च्या स्वयंसिद्धक~5 वरून
$\eq[(y'+z)'][0']$. स्वयंसिद्धक 1 वरून
$\eq[(y' +
    z)][0]$. पण व्याख्येने याचा अर्थ
$z < 0$; हे भाग~१ शी विसंगत आहे.
\end{proof}

\begin{explain}
संगणनक्षम फलनांचा संच म्हणजेच $\Th{Q}$ मध्ये निरूपणीय
फलनांचा संच, असे लक्षणन आपण दाखवले आहे. प्रत्यक्षात हा
पुरावा अधिक व्यापक आहे. निरूपणीयतेच्या व्याख्येवरून
$\Th{Q}$ चा विस्तार असलेल्या (किंवा ज्यात $\Th{Q}$ चे
अर्थनिर्धारण करता येते अशा) कोणत्याही उपपत्तीत
संगणनक्षम फलनांचे निरूपण करता येते, हे सहज दिसते.
उलट, सिद्धतेची संकल्पना संगणनक्षम असलेल्या कोणत्याही
!!{derivation} व्यवस्थेत प्रत्येक निरूपणीय फलन संगणनक्षम
असते, असे स्रोत म्हणतो. उदाहरणार्थ, संगणनक्षम फलनांचा संच
म्हणजे पेआनो अंकगणितात किंवा अगदी झर्मेलो--फ्रेंकेल संच
उपपत्तीत निरूपणीय फलनांचा संच, असे लक्षणन करता येते.
गोडेल यांनी नोंदवल्याप्रमाणे हे काहीसे आश्चर्यकारक आहे.
सिद्धतायोग्यतेचे प्रश्न कोणती उपपत्ती विचारात घेतो यावर
फार अवलंबून असतात, हे आपण पाहू: साधारणपणे स्वयंसिद्धके
जितकी बलवान, तितके अधिक सिद्ध करता येते. पण
स्वयंसिद्धकीय उपपत्तींच्या मोठ्या वर्गात निरूपणीय फलने
नेमकी संगणनक्षम फलनेच असतात, हा स्रोताचा दावा आहे.
%READERNOTE{OLINC-327: हा व्यापक उलट दावा सुसंगतता आणि भिन्न संख्यांकांची सिद्ध असमता यांशिवाय खोटा आहे; पूर्वीच्या विभागात OLINC-035 अंतर्गत अटी स्पष्ट केल्या आहेत.}
\end{explain}

\end{document}
